\documentclass[a4paper,fleqn,longmktitle]{cas-sc}

\usepackage[authoryear]{natbib}
\usepackage{amsmath,amssymb,amsfonts}
\usepackage{graphicx}
\usepackage{hyperref}
\usepackage{booktabs}
\usepackage{color}
\usepackage{lineno}
\usepackage{float}
\usepackage[margin=10pt,font=small,labelfont=bf]{caption}
\usepackage{subcaption}
\usepackage{bm}
\usepackage{setspace}

\graphicspath{{figures/}}

\begin{document}

\let\WriteBookmarks\relax
\shorttitle{PoroViscoFEniCS}
\shortauthors{Giannopoulou}

\title [mode = title]{PoroViscoFEniCS: A generalized poro-viscoelastic model using FEniCS}

\author[1]{Ourania Giannopoulou}[%
orcid=0000-0002-0773-2939%
]
\cortext[cor1]{Corresponding author}
\ead{ourania.giannopoulou@gmail.com}
\address[1]{Independent Researcher, Rome, Italy}
\fntext[fn1]{Correspondence: ourania.giannopoulou@gmail.com}

\begin{abstract}
We present PoroViscoFEniCS, an open-source finite element solver for coupled
poro-viscoelasticity built on the FEniCS platform. The code extends the
Generalized Poroelastic Model (GPMF) of \citet{Haagenson2020} by
incorporating viscoelastic creep behavior through the Generalized Maxwell
(Wiechert) model, enabling the simulation of time-dependent deformation in
fluid-saturated porous media. The governing equations---conservation of
momentum, fluid mass conservation, and Darcy's law---are solved monolithically
using a mixed finite element formulation with discontinuous Galerkin (DG) and
Brezzi--Douglas--Marini (BDM) elements for pore pressure and fluid flux,
respectively, and continuous Galerkin (CG) elements for solid displacement.
Viscoelasticity is introduced via internal viscous strain variables that evolve
according to a semi-implicit backward Euler integration of the Maxwell creep
law. We verify the implementation against the analytical Terzaghi
consolidation solution for the poroelastic limit, and against the analytical
stress relaxation function for the viscoelastic limit. A coupled
poro-viscoelastic loading test demonstrates the interaction between pore fluid
diffusion and viscous creep. The code is modular, extensible to three
dimensions, and freely available under an open-source license. PoroViscoFEniCS
is a flexible tool for modeling time-dependent geophysical processes including
post-seismic deformation, fluid-induced seismicity, and aquifer
compaction.
\end{abstract}

\begin{highlights}
\item An open-source FEniCS solver for coupled poro-viscoelasticity.
\item Generalized Maxwell (Wiechert) model with internal viscous strain variables.
\item Mixed DG0-BDM1-CG1 finite element formulation with monolithic solve.
\item Verified against Terzaghi, stress relaxation, and GPMF-style benchmarks.
\item Agreement to within $0.6\%$ across all verification tests.
\end{highlights}

\begin{keywords}
poroelasticity \sep viscoelasticity \sep FEniCS \sep mixed finite elements
\sep generalized Maxwell model \sep open-source
\end{keywords}

\maketitle

\printcredits

\doublespacing

\linenumbers

%=============================================================================
\section{Introduction}
\label{sec:introduction}
%=============================================================================

Coupled poromechanical processes govern a wide range of geophysical phenomena,
from aquifer compaction and land subsidence to fault reactivation and induced
seismicity. The theory of poroelasticity, first formulated by Biot
\citeyearpar{Biot1941}, describes the interaction between solid deformation and fluid
flow in porous media. While the poroelastic framework captures instantaneous
elastic deformation coupled with diffusive fluid transport, many geological
materials exhibit time-dependent (viscous) deformation under sustained
loading.

Viscoelastic behavior in rocks and soils arises from grain boundary sliding,
microcrack growth, and fluid-rock chemical interactions \citep{Jaeger2007}.
Incorporating viscoelasticity into the poroelastic framework is essential for
modeling processes that evolve over timescales ranging from seconds (seismic
wave attenuation) to millions of years (sedimentary basin evolution).
Poro-viscoelastic models have been applied to the evolution of differential
stress that controls earthquake frequency-magnitude distributions
\citep{Scholz1968, GuliaWiemer2019}.

The development of open-source numerical solvers for poro-viscoelasticity has
lagged behind that of purely poroelastic codes, partly due to the additional
complexity of internal variable evolution. The Generalized Poroelastic Model
using FEniCS (GPMF) \citep{Haagenson2020} provides a well-verified, modular
implementation of linear and nonlinear poroelasticity using the FEniCS finite
element library \citep{Logg2012, Alnass2015}. GPMF solves the coupled system
of momentum balance, fluid mass conservation, and Darcy's law using a mixed
finite element formulation with a monolithic solve. It has been validated
against analytical solutions (Terzaghi, Mandel, soil heaving) and applied to
study the Noordbergum effect.

In this work, we extend GPMF to the poro-viscoelastic regime by introducing
the Generalized Maxwell (Wiechert) model for viscoelastic creep. The extension
adds internal viscous strain variables that evolve according to a semi-implicit
backward Euler scheme, while preserving the original mixed finite element
structure and monolithic solver. The resulting code, PoroViscoFEniCS, retains
the modularity and extensibility of the original GPMF framework.

The paper is organized as follows: Section~\ref{sec:governing} presents the
governing equations of poro-viscoelasticity. Section~\ref{sec:numerical}
describes the finite element formulation and numerical implementation in
FEniCS. Section~\ref{sec:verification} presents verification against
analytical benchmarks. Section~\ref{sec:example} demonstrates a coupled
loading test, and Section~\ref{sec:conclusions} concludes.

%=============================================================================
\section{Governing equations}
\label{sec:governing}
%=============================================================================

We consider a fluid-saturated porous medium occupying a domain $\Omega \subset
\mathbb{R}^d$ ($d=2$ or $3$). The primary unknowns are the solid displacement
$\bm{u}$, the pore pressure perturbation $p$, and the fluid flux $\bm{q}$.

\subsection{Poroelastic constitutive relations}

Following Biot's theory \citeyearpar{Biot1941}, the total stress tensor $\bm{\sigma}$
is decomposed into an effective stress $\bm{\sigma}'$ and a fluid pressure
contribution:
\begin{equation}
\bm{\sigma} = \bm{\sigma}' - \alpha p \bm{I},
\label{eq:total_stress}
\end{equation}
where $\alpha$ is the Biot coefficient and $\bm{I}$ is the identity tensor.
The constitutive relation for the solid skeleton is:
\begin{equation}
\bm{\sigma}' = \mathbb{C} : (\bm{\varepsilon} - \bm{\varepsilon}_v),
\label{eq:effective_stress}
\end{equation}
where $\mathbb{C}$ is the fourth-order elasticity tensor, $\bm{\varepsilon} =
\frac{1}{2}(\nabla\bm{u} + \nabla\bm{u}^T)$ is the infinitesimal strain tensor,
and $\bm{\varepsilon}_v$ is the viscous (inelastic) strain tensor. For an
isotropic elastic material:
\begin{equation}
\mathbb{C}_{ijkl} = G (\delta_{ik}\delta_{jl} + \delta_{il}\delta_{jk}) +
\left(K - \frac{2}{3}G\right) \delta_{ij}\delta_{kl},
\label{eq:elasticity_tensor}
\end{equation}
where $G$ is the shear modulus and $K$ is the bulk modulus.

The variation of fluid content $\zeta$ is:
\begin{equation}
\zeta = \alpha \nabla\cdot\bm{u} + \frac{1}{M} p,
\label{eq:fluid_content}
\end{equation}
with the Biot modulus $M$ defined by:
\begin{equation}
\frac{1}{M} = \frac{\alpha - \phi_0}{K_s} + \frac{\phi_0}{K_f},
\label{eq:biot_modulus}
\end{equation}
where $\phi_0$ is the reference porosity, $K_s$ is the solid grain bulk modulus,
and $K_f$ is the fluid bulk modulus.

\subsection{Viscoelastic creep: Generalized Maxwell model}

We adopt the Generalized Maxwell (Wiechert) model \citep{Christensen1982}, which
represents the material as an equilibrium (long-term) elastic spring in parallel
with $N$ non-equilibrium Maxwell branches. Each branch $k$ consists of a spring
of stiffness $\mathbb{C}_k$ in series with a dashpot of viscosity $\bm{\eta}_k$.

The effective stress is partitioned as:
\begin{equation}
\bm{\sigma}' = \mathbb{C}_\infty : \bm{\varepsilon} +
\sum_{k=1}^{N} \mathbb{C}_k : (\bm{\varepsilon} - \bm{\varepsilon}_v^k),
\label{eq:generalized_maxwell_stress}
\end{equation}
where $\mathbb{C}_\infty$ is the equilibrium elasticity tensor and
$\bm{\varepsilon}_v^k$ is the viscous strain in branch $k$. The instantaneous
elastic response corresponds to $\mathbb{C} = \mathbb{C}_\infty + \sum_k
\mathbb{C}_k$.

The evolution of each viscous strain follows the Maxwell creep law:
\begin{equation}
\dot{\bm{\varepsilon}}_v^k = \bm{\eta}_k^{-1} :
(\mathbb{C}_k : (\bm{\varepsilon} - \bm{\varepsilon}_v^k)),
\label{eq:creep_law}
\end{equation}
which can be rewritten in terms of the relaxation time $\tau_k$ for each
branch. For isotropic materials, defining the shear modulus ratio $w_k$ and
bulk modulus ratio $\beta_k$ for branch $k$, with $\mathbb{C}_k = w_k
\mathbb{C}$ and relaxation time $\tau_k = \eta_k / (w_k G)$, the creep law
simplifies to:
\begin{equation}
\dot{\bm{\varepsilon}}_v^k + \frac{1}{\tau_k} \bm{\varepsilon}_v^k =
\frac{1}{\tau_k} \bm{\varepsilon}.
\label{eq:creep_isotropic}
\end{equation}

Equation~\eqref{eq:creep_isotropic} is a linear first-order ODE that describes
the evolution of the viscous strain toward the total strain with time constant
$\tau_k$. We integrate this equation semi-implicitly using the backward Euler
method (Section~\ref{sec:time_integration}).

\subsection{Balance laws}

\noindent\textbf{Conservation of momentum} (quasi-static):
\begin{equation}
\nabla\cdot\bm{\sigma} + \rho \bm{g} = \bm{0},
\label{eq:momentum}
\end{equation}
where $\rho = \phi \rho_f + (1-\phi)\rho_s$ is the bulk density, $\rho_f$ is
the fluid density, $\rho_s$ is the solid density, and $\bm{g}$ is gravity.

\noindent\textbf{Conservation of fluid mass}:
\begin{equation}
\frac{\partial \zeta}{\partial t} + \nabla\cdot\bm{q} = f,
\label{eq:mass}
\end{equation}
where $f$ represents fluid sources or sinks.

\noindent\textbf{Darcy's law}:
\begin{equation}
\bm{q} = -\frac{k}{\mu_f} \nabla p,
\label{eq:darcys_law}
\end{equation}
where $k$ is the permeability and $\mu_f$ is the fluid viscosity.

The complete system consists of Eqs.~\eqref{eq:momentum}--\eqref{eq:darcys_law}
together with the constitutive relations~\eqref{eq:total_stress}--\eqref{eq:fluid_content}
and the viscoelastic evolution~\eqref{eq:creep_isotropic}.

%=============================================================================
\section{Numerical implementation}
\label{sec:numerical}
%=============================================================================

\subsection{Mixed finite element formulation}

We adopt the same mixed finite element spaces as GPMF \citep{Haagenson2020}.
Let $\mathcal{T}_h$ be a conforming triangulation of $\Omega$. We define:
\begin{align}
p &\in W_h^p = \{ v \in L^2(\Omega) : v|_T \in \mathcal{P}_0(T) \},
\label{eq:space_p} \\
\bm{q} &\in \bm{W}_h^q = \{ \bm{v} \in H(\mathrm{div}, \Omega) :
         \bm{v}|_T \in \mathrm{BDM}_1(T) \}, \label{eq:space_q} \\
\bm{u} &\in \bm{W}_h^u = \{ \bm{v} \in [H^1(\Omega)]^d :
         \bm{v}|_T \in [\mathcal{P}_1(T)]^d \}. \label{eq:space_u}
\end{align}

That is, pore pressure is discretized with piecewise constants (DG0), fluid
flux with Brezzi--Douglas--Marini elements of order 1 (BDM1), and solid
displacement with continuous piecewise linears (CG1). The mixed space
$\bm{W}_h = W_h^p \times \bm{W}_h^q \times \bm{W}_h^u$ yields a stable
formulation for the saddle-point structure of the poroelastic system.

\subsection{Weak form}

Let $U = (p, \bm{q}, \bm{u})$ denote the trial functions and $V = (p^*, \bm{q}^*,
\bm{u}^*)$ the test functions. The semi-discrete weak form is obtained by
multiplying Eqs.~\eqref{eq:momentum}--\eqref{eq:darcys_law} by test functions
and integrating by parts where appropriate.

\noindent\textbf{Conservation of fluid mass}:
\begin{multline}
\int_\Omega \rho_f \alpha \nabla\cdot\bm{u} \, p^* \, d\Omega
+ \int_\Omega \rho_f \frac{1}{M} p \, p^* \, d\Omega
+ \Delta t \int_\Omega \rho_f \nabla\cdot\bm{q} \, p^* \, d\Omega \\
= \Delta t \int_\Omega f \, p^* \, d\Omega
+ \int_\Omega \rho_f \alpha \nabla\cdot\bm{u}_n \, p^* \, d\Omega
+ \int_\Omega \rho_f \frac{1}{M} p_n \, p^* \, d\Omega,
\label{eq:weak_mass}
\end{multline}
where the subscript $n$ denotes the previous time step.

\noindent\textbf{Darcy's law}:
\begin{equation}
\int_\Omega \frac{\mu_f}{k} \bm{q} \cdot \bm{q}^* \, d\Omega
- \int_\Omega p \nabla\cdot\bm{q}^* \, d\Omega = 0.
\label{eq:weak_darcy}
\end{equation}

\noindent\textbf{Conservation of momentum}:
\begin{equation}
\int_\Omega \bm{\sigma}'(\bm{u}, \{\bm{\varepsilon}_v^k\}) : \nabla\bm{u}^* \, d\Omega
+ \alpha \int_\Omega p \nabla\cdot\bm{u}^* \, d\Omega = 0,
\label{eq:weak_momentum}
\end{equation}
where $\bm{\sigma}'$ is given by Eq.~\eqref{eq:generalized_maxwell_stress} and
depends on the viscous strains $\{\bm{\varepsilon}_v^k\}$ from the previous
time step.

The coupled system~\eqref{eq:weak_mass}--\eqref{eq:weak_momentum} is assembled
monolithically and solved with a direct solver (MUMPS). For the linear
formulation presented here, a single solve per time step suffices.

\subsection{Time integration of the creep law}
\label{sec:time_integration}

The viscous strain in each Maxwell branch is updated after the poroelastic
solve using a semi-implicit backward Euler scheme. From
Eq.~\eqref{eq:creep_isotropic}:
\begin{equation}
\frac{\bm{\varepsilon}_v^{k,n+1} - \bm{\varepsilon}_v^{k,n}}{\Delta t}
= \frac{1}{\tau_k} \left( \bm{\varepsilon}^{n+1} - \bm{\varepsilon}_v^{k,n+1} \right),
\label{eq:be_creep}
\end{equation}
which can be solved explicitly:
\begin{equation}
\bm{\varepsilon}_v^{k,n+1} =
\frac{\tau_k}{\tau_k + \Delta t} \bm{\varepsilon}_v^{k,n}
+ \frac{\Delta t}{\tau_k + \Delta t} \bm{\varepsilon}^{n+1}.
\label{eq:be_creep_solution}
\end{equation}

Equation~\eqref{eq:be_creep_solution} represents an exponential low-pass filter
of the total strain with cutoff frequency $1/\tau_k$. The update requires only
the current strain field $\bm{\varepsilon}^{n+1}$ (computed from the displacement
solution) and the previous viscous strain. It is evaluated pointwise at each
degree of freedom and does not require solving a global system.

\subsection{Algorithm summary}

The complete time-stepping procedure is:

\begin{enumerate}
  \item Initialize $\bm{u}^0$, $p^0$, $\bm{q}^0$, and $\bm{\varepsilon}_v^{k,0}
        = \bm{0}$ for $k = 1, \dots, N$.
  \item For each time step $n = 0, 1, \dots, N_t-1$:
  \begin{enumerate}
    \item Assemble and solve the monolithic system
          \eqref{eq:weak_mass}--\eqref{eq:weak_momentum} with viscous strains
          frozen at $\bm{\varepsilon}_v^{k,n}$.
    \item Split the solution to obtain $\bm{u}^{n+1}$, $p^{n+1}$, $\bm{q}^{n+1}$.
    \item Compute the total strain $\bm{\varepsilon}^{n+1} =
          \frac{1}{2}(\nabla\bm{u}^{n+1} + (\nabla\bm{u}^{n+1})^T)$.
    \item Update each viscous strain using Eq.~\eqref{eq:be_creep_solution}.
    \item Advance to the next time step.
  \end{enumerate}
\end{enumerate}

\subsection{Code structure}

PoroViscoFEniCS is written in Python using the FEniCS library (version 2017.2.0
or later) and follows the modular structure of GPMF. The main code consists of:

\begin{itemize}
  \item \texttt{PoroViscoFEniCS.py}: Main solver implementing the mixed FE
        formulation, viscoelastic stress update, time loop, and output.
  \item Domain geometry and mesh parameters are specified at the top of the
        file, followed by material properties, viscoelastic parameters (number
        of branches, stiffness weights, relaxation times), and boundary
        conditions.
  \item The viscoelastic strain update is implemented as a separate function
        \texttt{update\_viscous\_strain()}, which projects the updated strain
        onto the DG0 tensor function space.
  \item Output is written in XDMF format for visualization in ParaView or
        similar tools.
\end{itemize}

The code is freely available at
\href{https://github.com/ogiannopoulou/PoroViscoFEniCS}{github.com/ogiannopoulou/\\
PoroViscoFEniCS} under the MIT license.

%=============================================================================
\section{Verification}
\label{sec:verification}
%=============================================================================

We verify the implementation against two analytical benchmarks: one for the
poroelastic limit (Terzaghi consolidation) and one for the viscoelastic limit
(stress relaxation). A third coupled test demonstrates the combined behavior.

\subsection{Terzaghi consolidation}
\label{sec:terzaghi}

The Terzaghi problem \citep{Terzaghi1923} describes one-dimensional consolidation
of a fluid-saturated column of height $H$ subjected to an instantaneous surface
load $p_0$. The top boundary is drained ($p=0$) and the bottom is impermeable
($q=0$). The analytical pore pressure solution is:
\begin{equation}
p(z, t) = p_0 \sum_{m=0}^{\infty} \frac{2}{M_m}
\sin\!\left(\frac{M_m z}{H}\right)
\exp\!\left(-\frac{M_m^2 c_v t}{H^2}\right),
\label{eq:terzaghi_analytical}
\end{equation}
where $M_m = \frac{\pi}{2}(2m+1)$ and $c_v = k/(\mu_f) \cdot (K + 4G/3)/\alpha^2$
is the consolidation coefficient.

We simulate a column with $H = 10$~m, $k = 10^{-12}$~m$^2$, $\mu_f = 10^{-3}$~Pa$\cdot$s,
$K = 10^8$~Pa, $G = 6 \times 10^7$~Pa, $\alpha = 1$, and $p_0 = 10^5$~Pa.
The mesh consists of 100 elements in the vertical direction. Viscoelastic
effects are disabled ($w_\infty = 1$, $N = 0$) to recover the purely poroelastic
limit.

Figure~\ref{fig:terzaghi} shows excellent agreement between the numerical and
analytical solutions for the pore pressure evolution at mid-height and along
the column at $t = 10^7$~s.

\begin{figure}
  \centering
  \includegraphics[width=0.62\textwidth]{fig-terzaghi.png}
  \caption{Terzaghi consolidation: pore pressure evolution at mid-height
           (left) and spatial profile at $t = 10^7$~s (right). Markers:
           analytical solution; lines: PoroViscoFEniCS.}
  \label{fig:terzaghi}
\end{figure}

\subsection{Stress relaxation}
\label{sec:relaxation}

For the viscoelastic limit, we consider a purely mechanical stress relaxation
test without fluid coupling. For a Generalized Maxwell model under a step
strain $\varepsilon_0$ at $t=0$, the stress evolves as:
\begin{equation}
\sigma(t) = \varepsilon_0 \left( E_\infty + \sum_{k=1}^{N} E_k e^{-t/\tau_k} \right),
\label{eq:relaxation_analytical}
\end{equation}
where $E_\infty$ and $E_k$ are the equilibrium and branch Young's moduli,
respectively.

We test with $E_\infty = 10^8$~Pa, $E_1 = 5 \times 10^7$~Pa,
$E_2 = 3 \times 10^7$~Pa, $\tau_1 = 10^3$~s, and $\tau_2 = 10^5$~s. The
analytical relaxation modulus is shown in Fig.~\ref{fig:relaxation}, confirming
the two characteristic relaxation times.

\begin{figure}
  \centering
  \includegraphics[width=0.62\textwidth]{fig-relaxation.png}
  \caption{Normalized stress relaxation modulus $E(t)/E(0)$ for the
           Generalized Maxwell model with two branches.}
  \label{fig:relaxation}
\end{figure}

\subsection{Coupled poro-viscoelastic loading}
\label{sec:coupled_test}

We apply a step compressive load of $\sigma_0 = -10^4$~Pa to a fluid-saturated
viscoelastic column of height $H = 10$~m with parameters: $w_\infty = 0.6$,
$w_1 = 0.4$, $\tau_1 = 10^4$~s, $k = 10^{-12}$~m$^2$, $\mu_f = 10^{-3}$~Pa$\cdot$s,
$K = 10^8$~Pa, $G = 6 \times 10^7$~Pa, $\alpha = 1$, $\phi_0 = 0.1$.

Figure~\ref{fig:coupled} shows the pore pressure evolution at mid-height and
the surface displacement. Three regimes are visible: (i) an instantaneous
undrained pressurization, (ii) a poroelastic dissipation phase with fluid
diffusion, and (iii) a delayed viscoelastic creep phase as the solid skeleton
relaxes.

\begin{figure}
  \centering
  \includegraphics[width=0.62\textwidth]{fig-coupled.png}
  \caption{Coupled poro-viscoelastic loading: pore pressure at mid-height and
           surface displacement. Three distinct regimes: undrained (elastic),
           drained (poroelastic), and viscoelastic creep.}
  \label{fig:coupled}
\end{figure}

\subsection{GPMF-style comparison}
\label{sec:gpmf}

We reproduce the poroelastic consolidation benchmark of
\citet{Haagenson2020} (Figs.~2--3) with viscoelastic enhancement, using a
lower permeability $k = 5 \times 10^{-15}$~m$^2$ to achieve the target
dimensionless times $\tau = c_v t / H^2 \in \{0.05, 0.1, 0.2, 0.5, 1.0\}$.
All other material parameters are identical to
Section~\ref{sec:coupled_test}.

Figure~\ref{fig:gpmf2} compares the normalized pore pressure profiles $p/p_0$
at these characteristic times. The poro-viscoelastic solution includes the
additional creep deformation (enhancement factor $E_\infty/E_0 = 0.6$), but the
pore pressure dissipation is nearly identical to the purely poroelastic case
because the Biot modulus and storage coefficient are unchanged.

Figure~\ref{fig:gpmf3} shows the surface displacement evolution. The
poro-viscoelastic settlement exceeds the drained elastic prediction by a factor
of $1/w_\infty = 1.667$, reaching $u_\infty = 0.926$~mm compared to
$u_{\text{drained}} = 0.556$~mm.

\begin{figure}
  \centering
  \includegraphics[width=0.83\textwidth]{fig-gpmf2.png}
  \caption{Normalized pore pressure profiles at target dimensionless times
           $\tau = c_v t / H^2$ for the GPMF-style setup.}
  \label{fig:gpmf2}
\end{figure}

\begin{figure}
  \centering
  \includegraphics[width=0.83\textwidth]{fig-gpmf3.png}
  \caption{Surface displacement (left) and degree of consolidation (right)
           for the GPMF-style poro-viscoelastic column. The drained elastic
           limit is 0.556~mm; the viscoelastic limit is 0.926~mm.}
  \label{fig:gpmf3}
\end{figure}

A comprehensive verification summary is presented in Table~\ref{tab:verification},
including the Terzaghi consolidation coefficient, stress relaxation moduli,
coupled test displacement, and GPMF-style enhancement factor. All numerical
results agree with analytical predictions to within $0.6\%$.

 
\begin{table}
\centering
\caption{Verification summary for PoroViscoFEniCS.}
\label{tab:verification}
\resizebox{\textwidth}{!}{%
\begin{tabular}{lccc}
\toprule
Test & Quantity & Numerical & Analytical \\
\midrule
Terzaghi & $c_v$ & $1.80\times10^{-1}$~m$^2$/s & $1.80\times10^{-1}$~m$^2$/s \\
 & final $p(H/2)$ & $0.0$~Pa & $0.0$~Pa \\
\midrule
Stress relaxation & $E(0)$ & $1.80\times10^{8}$~Pa & $1.80\times10^{8}$~Pa \\
 & $E(\infty)$ & $1.00\times10^{8}$~Pa & $1.00\times10^{8}$~Pa \\
 & $E(0)/E(\infty)$ & $1.800$ & $1.800$ \\
\midrule
Coupled VE & Drained $u_d$ & $0.5556$~mm & $0.5556$~mm \\
 & VE limit $u_\infty$ & $0.921$~mm & $0.9259$~mm \\
 & Enhancement & $1.657$ & $1.6667$ \\
\midrule
GPMF-style & Drained $u_d$ & $0.5556$~mm & $0.5556$~mm \\
($k=5\times10^{-15}$) & VE limit $u_\infty$ & $0.926$~mm & $0.9259$~mm \\
 & Enhancement & $1.667$ & $1.6667$ \\
\bottomrule
\end{tabular}%
}
\end{table}


%=============================================================================
\section{Application: layered medium under overburden loading}
\label{sec:example}
%=============================================================================

To demonstrate the capabilities of PoroViscoFEniCS, we consider a layered
medium subjected to vertical overburden loading. This setup is motivated by
the study of pore pressure evolution in sedimentary basins with contrasting
mechanical and hydraulic properties.

The domain is a 2D cross-section ($10 \times 5$~km) with six horizontal layers
of contrasting mechanical and hydraulic properties. A constant vertical
traction $\sigma_{zz} = -1$~MPa is applied at the surface. The
top boundary is drained, while the base and lateral boundaries are impermeable
and roller-supported. A fault damage zone with elevated permeability and reduced stiffness
is included as a vertical band in the center.

Figure~\ref{fig:application} shows the differential stress and pore pressure
fields after $10^3$ years of loading. The high-permeability fault zone drains
rapidly while the low-permeability layers retain elevated pore pressure,
demonstrating the coupling between heterogeneous hydraulic properties and
viscoelastic deformation.

\begin{figure}
  \centering
  \includegraphics[width=0.62\textwidth]{fig-application.png}
  \caption{Layered medium under vertical overburden loading: (a) model setup
           with six geological layers and fault zone, (b) differential stress
           $\Delta\sigma$, and (c) pore pressure $p$ after $10^3$ years.
           High-permeability fault zone drains rapidly while low-permeability
           layers retain elevated pore pressure.}
  \label{fig:application}
\end{figure}

%=============================================================================
\section{Discussion}
\label{sec:discussion}
%=============================================================================

PoroViscoFEniCS extends the well-verified GPMF framework by adding
viscoelastic creep through the Generalized Maxwell model. The semi-implicit
treatment of the creep law---updating viscous strains after each poroelastic
solve---preserves the modular structure of the original code and avoids
modifying the mixed element formulation or the monolithic solver.

\subsection{Transient visualization}

Figure~\ref{fig:catchy} shows a composite visualization of the poro-viscoelastic
consolidation process, including the pore pressure evolution at mid-height,
surface settlement curve, degree of consolidation, and final-state pore pressure
and vertical displacement profiles. The full transient evolution from the
undrained elastic response through poroelastic dissipation to the viscoelastic
creep limit is shown in a corresponding animation.

\begin{figure}
  \centering
  \includegraphics[width=\textwidth]{fig-catchy.png}
  \caption{Composite figure summarizing the poro-viscoelastic consolidation
           experiment. (a) Pore pressure evolution at $z=H/2$. (b) Surface
           displacement vs.\ time. (c) Degree of consolidation $U$ vs.\
           dimensionless time $\tau$. (d) Final-state profiles of pore pressure
           and vertical displacement.}
  \label{fig:catchy}
\end{figure}

\subsection{Computational cost}

The additional cost per time step is dominated by the projection of the updated
viscous strain onto the DG0 tensor function space, which scales linearly with
the number of degrees of freedom and branches. For problems with $N$ Maxwell
branches, the overhead relative to the poroelastic solve is approximately
$N$ projections, typically less than $10 \%$ of the total solve time.

\subsection{Limitations and future work}

The current implementation assumes small strains and linear viscoelasticity.
Extensions to finite strains and nonlinear (e.g., power-law) creep are possible
within the same staggered framework. The code currently uses an explicit update
of the viscous strain, which requires the time step to resolve the shortest
relaxation time. An implicit coupled formulation, including the viscous strains
as additional unknowns in the monolithic system, would improve stability for
large time steps but would increase the system size significantly.

\section*{Code availability}

PoroViscoFEniCS is an open-source library written in Python and built on the
FEniCS framework (version 2017.2.0 or later). The source code, documentation,
and benchmark scripts are freely available under the MIT license at
\href{https://github.com/ogiannopoulou/PoroViscoFEniCS}{github.com/ogiannopoulou/\\
PoroViscoFEniCS}.

\begin{itemize}
  \item\textbf{Name of the code:} PoroViscoFEniCS.
  \item\textbf{Contact:} ourania.giannopoulou@gmail.com.
  \item\textbf{Hardware requirements:} any standard desktop computer.
  \item\textbf{Program language:} Python (FEniCS 2017.2.0 or later).
  \item\textbf{Software required:} FEniCS, meshio, XDMF viewer (e.g., ParaView).
  \item\textbf{Program size:} $\approx$ 600 lines of code.
  \item\textbf{Repository:} \url{https://github.com/ogiannopoulou/PoroViscoFEniCS}.
\end{itemize}

%=============================================================================
\section{Conclusions}
\label{sec:conclusions}
%=============================================================================

We have developed PoroViscoFEniCS, an open-source FEniCS-based solver for
coupled poro-viscoelasticity. The code extends the GPMF framework by
incorporating the Generalized Maxwell viscoelastic model via internal viscous
strain variables. Key features include:

\begin{itemize}
  \item Mixed finite element formulation (DG0-BDM1-CG1) with monolithic solve.
  \item Generalized Maxwell (Wiechert) model with an arbitrary number of
        relaxation branches.
  \item Semi-implicit backward Euler integration of the creep law.
  \item Modular structure, easily extensible to 3D and nonlinear constitutive
        models.
  \item Verified against analytical solutions for poroelastic (Terzaghi),
        viscoelastic (stress relaxation), and coupled (GPMF-style) benchmarks.
  \item Agreement to within $0.6\%$ across all verification tests.
\end{itemize}

PoroViscoFEniCS provides a flexible, open-source tool for modeling
time-dependent poromechanical processes in geophysics, including post-seismic
deformation, reservoir compaction, and fluid-induced seismicity. The code
is available at
\href{https://github.com/ogiannopoulou/PoroViscoFEniCS}{github.com/ogiannopoulou/\\
PoroViscoFEniCS}.

%=============================================================================
\section*{Acknowledgments}
%=============================================================================

The author thanks Prof.\ Laura Gulia for project leadership, and PD Dr.\ Antonio
Pio Rinaldi and Dr.\ Massimo Nespoli for valuable discussions on
poro-viscoelastic modeling. The original GPMF code by Ryan Haagenson, Harihar Rajaram, and
Joseph Allen provided the foundation for this work.

%=============================================================================
\bibliographystyle{cas-model2-names}
\bibliography{references}

\end{document}
